\documentclass[12pt,a4paper]{article}

% ── Pacotes ──────────────────────────────────────────────────────────────────
\usepackage[utf8]{inputenc}
\usepackage[T1]{fontenc}
\usepackage[portuguese]{babel}
\usepackage[margin=2.5cm]{geometry}
\usepackage{xcolor}
\usepackage{titlesec}
\usepackage{titletoc}
\usepackage{tocloft}
\usepackage{hyperref}
\usepackage{fontawesome5}
\usepackage{booktabs}
\usepackage{longtable}
\usepackage{array}
\usepackage{enumitem}
\usepackage{fancyhdr}
\usepackage{graphicx}
\usepackage{mdframed}
\usepackage{lmodern}
\usepackage{microtype}
\usepackage{parskip}
\usepackage{tcolorbox}
\tcbuselibrary{skins, breakable}

% ── Cores ────────────────────────────────────────────────────────────────────
\definecolor{orbisblue}{HTML}{4F46E5}
\definecolor{orbisDark}{HTML}{1E1B4B}
\definecolor{orbisLight}{HTML}{EEF2FF}
\definecolor{orbisGreen}{HTML}{059669}
\definecolor{orbisRed}{HTML}{DC2626}
\definecolor{orbisYellow}{HTML}{D97706}
\definecolor{orbisMid}{HTML}{6366F1}
\definecolor{codebg}{HTML}{F1F5F9}
\definecolor{tabhead}{HTML}{312E81}

% ── Hiperligações ─────────────────────────────────────────────────────────────
\hypersetup{
  colorlinks=true,
  linkcolor=orbisblue,
  urlcolor=orbisMid,
  pdftitle={Orbis — Manual do Utilizador},
  pdfauthor={Orion Technologies},
}

% ── Cabeçalho / Rodapé ───────────────────────────────────────────────────────
\pagestyle{fancy}
\fancyhf{}
\fancyhead[L]{\textcolor{orbisblue}{\textbf{Orbis}}}
\fancyhead[R]{\textcolor{gray}{\small Manual do Utilizador}}
\fancyfoot[C]{\textcolor{gray}{\small \thepage}}
\renewcommand{\headrulewidth}{0.4pt}
\renewcommand{\headrule}{\hbox to\headwidth{\color{orbisblue}\leaders\hrule height \headrulewidth\hfill}}

% ── Títulos ───────────────────────────────────────────────────────────────────
\titleformat{\section}
  {\Large\bfseries\color{orbisDark}}
  {\textcolor{orbisblue}{\thesection.}}{0.6em}{}
  [\textcolor{orbisblue}{\titlerule[1.2pt]}]

\titleformat{\subsection}
  {\large\bfseries\color{orbisblue}}
  {\thesubsection.}{0.5em}{}

\titleformat{\subsubsection}
  {\normalsize\bfseries\color{orbisMid}}
  {}{0em}{}

% ── Caixas personalizadas ─────────────────────────────────────────────────────
\tcbset{
  cmdbox/.style={
    enhanced, breakable,
    colback=codebg, colframe=orbisblue,
    fonttitle=\bfseries\color{white},
    coltitle=white, attach boxed title to top left,
    boxed title style={colback=orbisblue, rounded corners},
    arc=4pt, boxrule=0.6pt,
    left=6pt, right=6pt, top=6pt, bottom=6pt,
  },
  notebox/.style={
    enhanced, breakable,
    colback=orbisLight, colframe=orbisMid,
    arc=4pt, boxrule=0.8pt,
    left=8pt, right=8pt, top=6pt, bottom=6pt,
  },
  warnbox/.style={
    enhanced,
    colback=yellow!8, colframe=orbisYellow,
    arc=4pt, boxrule=0.8pt,
    left=8pt, right=8pt, top=6pt, bottom=6pt,
  },
  dangerbox/.style={
    enhanced,
    colback=red!6, colframe=orbisRed,
    arc=4pt, boxrule=0.8pt,
    left=8pt, right=8pt, top=6pt, bottom=6pt,
  },
}

% ── Comando inline para comandos ──────────────────────────────────────────────
\newcommand{\cmd}[1]{\texttt{\textcolor{orbisblue}{\textbf{#1}}}}
\newcommand{\param}[1]{\texttt{\textcolor{orbisMid}{#1}}}

% ── Tabela de comandos ────────────────────────────────────────────────────────
\newcolumntype{C}[1]{>{\centering\arraybackslash}p{#1}}
\newcolumntype{L}[1]{>{\raggedright\arraybackslash}p{#1}}

% ─────────────────────────────────────────────────────────────────────────────
\begin{document}

% ── CAPA ─────────────────────────────────────────────────────────────────────
\begin{titlepage}
  \pagecolor{orbisDark}
  \color{white}
  \centering
  \vspace*{3cm}

  {\fontsize{60}{60}\selectfont \textcolor{orbisMid}{\faRobot}}

  \vspace{1cm}
  {\fontsize{48}{48}\selectfont \textbf{Orbis}}

  \vspace{0.4cm}
  {\large \textcolor{orbisMid}{Bot Inteligente de Gestão de Grupos WhatsApp}}

  \vspace{2cm}
  \textcolor{orbisMid}{\rule{10cm}{0.6pt}}
  \vspace{0.6cm}

  {\Large \textbf{Manual do Utilizador}}\\[0.4cm]
  {\normalsize Versão 1.0}

  \vfill

  \textcolor{orbisMid}{\rule{10cm}{0.4pt}}\\[0.5cm]
  {\small Desenvolvido por \textbf{Orion Technologies}}\\[0.2cm]
  {\small \textcolor{gray}{\url{https://oriontechn.netlify.app}}}\\[0.6cm]
  {\small \textcolor{gray}{\today}}
\end{titlepage}
\pagecolor{white}\color{black}

% ── ÍNDICE ────────────────────────────────────────────────────────────────────
\newpage
\tableofcontents
\newpage


% ═════════════════════════════════════════════════════════════════════════════
\section{Introdução}
% ═════════════════════════════════════════════════════════════════════════════

O \textbf{Orbis} é um bot autónomo de gestão de grupos WhatsApp com inteligência artificial integrada. Modera grupos automaticamente, responde a comandos, conversa com os membros e notifica administradores — tudo sem intervenção manual.

\begin{tcolorbox}[notebox]
\textbf{\textcolor{orbisblue}{\faInfoCircle\ O que o Orbis faz por ti:}}
\begin{itemize}[leftmargin=1.2em, itemsep=2pt]
  \item Detecta e remove spam, links proibidos e flood automaticamente
  \item Aplica punições progressivas: aviso $\to$ mute $\to$ expulsão $\to$ ban
  \item Responde a perguntas e conversa com os membros via IA (Groq)
  \item Executa comandos de moderação por linguagem natural (admins)
  \item Envia resumos diários da conversa do grupo
  \item Gere agendamentos recorrentes de mensagens
  \item Mantém um sistema de reputação e níveis para os membros
\end{itemize}
\end{tcolorbox}

\subsection{Tecnologias}

\begin{center}
\begin{tabular}{L{3.5cm} L{9cm}}
\toprule
\rowcolor{tabhead}\textcolor{white}{\textbf{Tecnologia}} & \textcolor{white}{\textbf{Função}} \\
\midrule
Baileys v7 & Conexão com o WhatsApp \\
\rowcolor{orbisLight} Supabase (PostgreSQL) & Base de dados persistente \\
Groq & Motor de Inteligência Artificial \\
\rowcolor{orbisLight} Express & API REST interna \\
Node.js v18+ & Ambiente de execução \\
\bottomrule
\end{tabular}
\end{center}

% ═════════════════════════════════════════════════════════════════════════════
\section{Níveis de Acesso}
% ═════════════════════════════════════════════════════════════════════════════

O Orbis distingue três perfis de utilizador dentro de cada grupo:

\begin{center}
\begin{tabular}{L{2.8cm} C{1.8cm} L{8cm}}
\toprule
\rowcolor{tabhead}\textcolor{white}{\textbf{Perfil}} & \textcolor{white}{\textbf{Símbolo}} & \textcolor{white}{\textbf{Capacidades}} \\
\midrule
\textbf{Membro} & \faUser & Comandos gerais, sondagens, denúncias, reputação, resumos \\
\rowcolor{orbisLight}\textbf{Administrador} & \faCrown & Tudo do membro + moderação, configuração, agendamentos, anúncios \\
\textbf{Owner} & \faShieldAlt & Acesso total; imune a todas as acções de moderação \\
\bottomrule
\end{tabular}
\end{center}

\begin{tcolorbox}[warnbox]
\textbf{\textcolor{orbisYellow}{\faExclamationTriangle\ Importante:}} Os administradores do grupo no WhatsApp são automaticamente reconhecidos como \textbf{Administradores} pelo Orbis. Não é necessária configuração adicional.
\end{tcolorbox}


% ═════════════════════════════════════════════════════════════════════════════
\section{Comandos Gerais}
\label{sec:geral}
% ═════════════════════════════════════════════════════════════════════════════

Estes comandos estão disponíveis para \textbf{todos os membros} do grupo.

\subsection{Informação e Navegação}

\begin{tcolorbox}[cmdbox, title={\faList\ /menu}]
Mostra a lista resumida de todos os comandos disponíveis.\\
\textbf{Uso:} \cmd{/menu}
\end{tcolorbox}

\begin{tcolorbox}[cmdbox, title={\faInfoCircle\ /info}]
Mostra informações do grupo: nome, número de membros, admins e data de criação.\\
\textbf{Uso:} \cmd{/info}
\end{tcolorbox}

\begin{tcolorbox}[cmdbox, title={\faCrown\ /admins}]
Lista todos os administradores do grupo com menção.\\
\textbf{Uso:} \cmd{/admins}
\end{tcolorbox}

\begin{tcolorbox}[cmdbox, title={\faScroll\ /regras}]
Mostra as regras do grupo. Se não foram definidas, apresenta as regras padrão.\\
\textbf{Uso:} \cmd{/regras}
\end{tcolorbox}

\begin{tcolorbox}[cmdbox, title={\faChartBar\ /stats}]
Estatísticas do grupo: total de membros, admins e membros regulares.\\
\textbf{Uso:} \cmd{/stats}
\end{tcolorbox}

\begin{tcolorbox}[cmdbox, title={\faBook\ /instrucoes}]
Instruções detalhadas organizadas por secção.\\
\textbf{Uso:} \cmd{/instrucoes} \param{[secção]}\\
\textbf{Secções disponíveis:} \param{geral} \textbar\ \param{moderacao} \textbar\ \param{config} \textbar\ \param{ia} \textbar\ \param{extra}\\
\textbf{Exemplo:} \cmd{/instrucoes moderacao}
\end{tcolorbox}

\subsection{Interacção e Participação}

\begin{tcolorbox}[cmdbox, title={\faPoll\ /sondagem}]
Cria uma enquete nativa do WhatsApp com 2 a 12 opções.\\
\textbf{Uso:} \cmd{/sondagem} \param{<pergunta> | opção1 | opção2 ...}\\
\textbf{Exemplo:} \cmd{/sondagem Melhor dia para reunião? | Segunda | Quarta | Sexta}
\end{tcolorbox}

\begin{tcolorbox}[cmdbox, title={\faFlag\ /denunciar}]
Denuncia uma mensagem aos administradores. Deves \textbf{responder} à mensagem que queres denunciar antes de usar o comando.\\
\textbf{Uso:} Responde a uma mensagem $\to$ \cmd{/denunciar}
\end{tcolorbox}

\begin{tcolorbox}[cmdbox, title={\faStar\ /reputacao}]
Mostra o nível e pontos de reputação de um membro (ou os teus, se não mencionares ninguém).\\
\textbf{Uso:} \cmd{/reputacao} \param{[@membro]}\\
\textbf{Exemplo:} \cmd{/reputacao @João}
\end{tcolorbox}

\begin{tcolorbox}[cmdbox, title={\faTrophy\ /top}]
Mostra o top 10 de membros com mais reputação no grupo.\\
\textbf{Uso:} \cmd{/top}
\end{tcolorbox}

\begin{tcolorbox}[cmdbox, title={\faAlignLeft\ /resumo}]
A Orbis resume as últimas N mensagens da conversa usando IA.\\
\textbf{Uso:} \cmd{/resumo} \param{[N]}\\
\textbf{Padrão:} 20 mensagens. \textbf{Exemplo:} \cmd{/resumo 50}
\end{tcolorbox}


% ═════════════════════════════════════════════════════════════════════════════
\section{Comandos de Moderação}
\label{sec:moderacao}
% ═════════════════════════════════════════════════════════════════════════════

Estes comandos são exclusivos para \textbf{administradores} do grupo.

\subsection{Avisos}

\begin{tcolorbox}[cmdbox, title={\faExclamationTriangle\ /warn}]
Dá um aviso a um membro. Ao atingir o limite configurado, o membro é expulso automaticamente.\\
\textbf{Uso:} \cmd{/warn} \param{@membro [motivo]}\\
\textbf{Exemplo:} \cmd{/warn @Pedro spam repetido}
\end{tcolorbox}

\begin{tcolorbox}[cmdbox, title={\faEyeSlash\ /warnings}]
Mostra o número de avisos actuais de um membro.\\
\textbf{Uso:} \cmd{/warnings} \param{@membro}
\end{tcolorbox}

\begin{tcolorbox}[cmdbox, title={\faBroom\ /clearwarn}]
Apaga todos os avisos de um membro, dando-lhe uma nova oportunidade.\\
\textbf{Uso:} \cmd{/clearwarn} \param{@membro}
\end{tcolorbox}

\begin{tcolorbox}[cmdbox, title={\faHistory\ /historico}]
Mostra o histórico completo de um membro: avisos, bans e denúncias recebidas.\\
\textbf{Uso:} \cmd{/historico} \param{@membro}
\end{tcolorbox}

\subsection{Expulsão e Ban}

\begin{tcolorbox}[cmdbox, title={\faUserMinus\ /kick}]
Expulsa o membro do grupo imediatamente. O membro pode ser readmitido manualmente.\\
\textbf{Uso:} \cmd{/kick} \param{@membro}
\end{tcolorbox}

\begin{tcolorbox}[cmdbox, title={\faHammer\ /ban}]
Bane o membro: expulsa-o e regista-o na lista de banidos, impedindo readmissão automática.\\
\textbf{Uso:} \cmd{/ban} \param{@membro [motivo]}\\
\textbf{Exemplo:} \cmd{/ban @Rui burla e fraude}
\end{tcolorbox}

\begin{tcolorbox}[cmdbox, title={\faList\ /banidos}]
Lista os últimos 20 membros banidos do grupo com o motivo do ban.\\
\textbf{Uso:} \cmd{/banidos}
\end{tcolorbox}

\begin{tcolorbox}[cmdbox, title={\faUserPlus\ /adicionar}]
Readmite um membro banido pelo número de telefone e remove-o da lista de banidos.\\
\textbf{Uso:} \cmd{/adicionar} \param{<número>}\\
\textbf{Exemplo:} \cmd{/adicionar 244912345678}
\end{tcolorbox}

\subsection{Mute e Controlo de Voz}

\begin{tcolorbox}[cmdbox, title={\faVolumeMute\ /mute}]
Muta um membro (remove permissão de envio). Com tempo definido, o mute expira automaticamente.\\
\textbf{Uso:} \cmd{/mute} \param{@membro [tempo]}\\
\textbf{Formatos de tempo:} \param{10m} (minutos), \param{2h} (horas), \param{1d} (dias)\\
\textbf{Exemplo:} \cmd{/mute @Ana 30m}
\end{tcolorbox}

\begin{tcolorbox}[cmdbox, title={\faVolumeUp\ /unmute}]
Remove o mute de um membro antes do tempo expirar.\\
\textbf{Uso:} \cmd{/unmute} \param{@membro}
\end{tcolorbox}

\subsection{Gestão de Administradores}

\begin{tcolorbox}[cmdbox, title={\faCrown\ /promover \ e \ /despromover}]
Promove ou despromove um membro a administrador do grupo.\\
\textbf{Uso:} \cmd{/promover} \param{@membro} \quad \cmd{/despromover} \param{@membro}
\end{tcolorbox}

\subsection{Controlo do Grupo}

\begin{tcolorbox}[cmdbox, title={\faLock\ /trancar \ e \ \faLockOpen\ /abrir}]
\textbf{/trancar} — Só administradores podem enviar mensagens.\\
\textbf{/abrir} — Todos os membros podem enviar mensagens.\\
\textbf{Uso:} \cmd{/trancar} \quad \cmd{/abrir}
\end{tcolorbox}

\begin{tcolorbox}[cmdbox, title={\faBullhorn\ /anunciar}]
Envia uma mensagem formatada como anúncio oficial da administração.\\
\textbf{Uso:} \cmd{/anunciar} \param{<texto>}\\
\textbf{Exemplo:} \cmd{/anunciar Reunião amanhã às 18h. Presença obrigatória.}
\end{tcolorbox}

\begin{tcolorbox}[cmdbox, title={\faTrash\ /deletar}]
Apaga uma mensagem específica. Deves \textbf{responder} à mensagem que queres apagar.\\
\textbf{Uso:} Responde a uma mensagem $\to$ \cmd{/deletar}\\
\textit{Nota: O bot precisa de ser administrador do grupo para apagar mensagens de outros.}
\end{tcolorbox}

\subsection{Agendamentos}

\begin{tcolorbox}[cmdbox, title={\faClock\ /agendamento}]
Agenda uma mensagem única para ser enviada após um determinado tempo.\\
\textbf{Uso:} \cmd{/agendamento} \param{<tempo> <mensagem>}\\
\textbf{Exemplo:} \cmd{/agendamento 1h Reunião em breve!}
\end{tcolorbox}

\begin{tcolorbox}[cmdbox, title={\faClock\ /agendarcorrente}]
Cria um agendamento recorrente — diário ou apenas dias úteis.\\
\textbf{Uso:} \cmd{/agendarcorrente} \param{<DAILY|MON-FRI> <HH:MM> <mensagem>}\\
\textbf{Exemplos:}
\begin{itemize}[leftmargin=1.2em, itemsep=1pt]
  \item \cmd{/agendarcorrente DAILY 09:00 Bom dia a todos!}
  \item \cmd{/agendarcorrente MON-FRI 08:30 Reunião em 30 minutos!}
\end{itemize}
\end{tcolorbox}

\begin{tcolorbox}[cmdbox, title={\faListAlt\ /agendamentos}]
Lista todos os agendamentos recorrentes activos no grupo.\\
\textbf{Uso:} \cmd{/agendamentos}
\end{tcolorbox}

\begin{tcolorbox}[cmdbox, title={\faTimes\ /cancelaragendamento}]
Cancela um agendamento recorrente pelo seu ID (visível em \cmd{/agendamentos}).\\
\textbf{Uso:} \cmd{/cancelaragendamento} \param{<id>}
\end{tcolorbox}

\begin{tcolorbox}[cmdbox, title={\faExclamationCircle\ /denuncias}]
Mostra as últimas denúncias registadas no grupo.\\
\textbf{Uso:} \cmd{/denuncias}
\end{tcolorbox}


% ═════════════════════════════════════════════════════════════════════════════
\section{Configuração do Grupo}
\label{sec:config}
% ═════════════════════════════════════════════════════════════════════════════

\subsection{Ver e Alterar Configurações}

\begin{tcolorbox}[cmdbox, title={\faCog\ /config}]
Sem argumentos, mostra todas as configurações actuais do grupo.\\
Com argumentos, altera uma configuração específica.\\
\textbf{Uso:} \cmd{/config} \param{[opção] [valor]}
\end{tcolorbox}

\subsubsection{Opções disponíveis do /config}

\begin{center}
\begin{longtable}{L{4.5cm} L{4cm} L{5cm}}
\toprule
\rowcolor{tabhead}\textcolor{white}{\textbf{Comando}} & \textcolor{white}{\textbf{Valores}} & \textcolor{white}{\textbf{Descrição}} \\
\midrule
\endhead
\cmd{/config spam on|off} & \param{on} / \param{off} & Activa/desactiva anti-spam \\
\rowcolor{orbisLight}\cmd{/config link on|off} & \param{on} / \param{off} & Activa/desactiva anti-link \\
\cmd{/config flood on|off} & \param{on} / \param{off} & Activa/desactiva anti-flood \\
\rowcolor{orbisLight}\cmd{/config welcome on|off} & \param{on} / \param{off} & Activa/desactiva boas-vindas \\
\cmd{/config warnings N} & número inteiro & Define o limite de avisos antes da expulsão \\
\rowcolor{orbisLight}\cmd{/config reputacao on|off} & \param{on} / \param{off} & Activa/desactiva sistema de reputação \\
\cmd{/config silencio H1 H2} & horas (0--23) & Modo silêncio das H1h às H2h \\
\rowcolor{orbisLight}\cmd{/config whitelist add D} & domínio & Adiciona domínio à whitelist de links \\
\cmd{/config whitelist remove D} & domínio & Remove domínio da whitelist \\
\rowcolor{orbisLight}\cmd{/config whitelist clear} & --- & Limpa toda a whitelist \\
\bottomrule
\end{longtable}
\end{center}

\subsection{Personalização}

\begin{tcolorbox}[cmdbox, title={\faEdit\ /setregras}]
Define as regras personalizadas do grupo (substitui as regras padrão).\\
\textbf{Uso:} \cmd{/setregras} \param{<texto das regras>}
\end{tcolorbox}

\begin{tcolorbox}[cmdbox, title={\faHandshake\ /setwelcome}]
Define a mensagem de boas-vindas para novos membros.\\
\textbf{Variáveis disponíveis:} \param{\{nome\}} (nome do membro) e \param{\{grupo\}} (nome do grupo)\\
\textbf{Uso:} \cmd{/setwelcome} \param{<mensagem>}\\
\textbf{Exemplo:} \cmd{/setwelcome Olá \{nome\}, bem-vindo(a) ao \{grupo\}! Lê as /regras.}
\end{tcolorbox}

\begin{tcolorbox}[cmdbox, title={\faBell\ /setadmingroup}]
Configura um grupo dedicado para receber notificações de moderação (bans, kicks, avisos, denúncias).\\
\textbf{Uso:} No grupo de admins, envia \cmd{/setadmingroup} \param{<ID do grupo principal>}\\
\textbf{Exemplo:} \cmd{/setadmingroup 120363427347409993@g.us}\\[4pt]
\textit{O ID do grupo aparece no terminal quando envias uma mensagem nesse grupo.}
\end{tcolorbox}

\subsection{Atalhos de Moderação}

\begin{center}
\begin{tabular}{L{5cm} L{8cm}}
\toprule
\rowcolor{tabhead}\textcolor{white}{\textbf{Comando}} & \textcolor{white}{\textbf{Equivalente a}} \\
\midrule
\cmd{/antilink on|off} & \cmd{/config link on|off} \\
\rowcolor{orbisLight}\cmd{/antispam on|off} & \cmd{/config spam on|off} \\
\cmd{/antiflood on|off} & \cmd{/config flood on|off} \\
\bottomrule
\end{tabular}
\end{center}

% ═════════════════════════════════════════════════════════════════════════════
\section{Inteligência Artificial}
\label{sec:ia}
% ═════════════════════════════════════════════════════════════════════════════

\subsection{Como Activar a IA}

A Orbis responde automaticamente quando:

\begin{itemize}[leftmargin=1.5em, itemsep=3pt]
  \item Mencionas \textbf{@Orbis} na mensagem
  \item Escreves a palavra \textbf{``orbis''} no texto
  \item Respondes directamente a uma mensagem da Orbis
  \item Envias uma mensagem \textbf{privada} directamente ao bot
\end{itemize}

\begin{tcolorbox}[notebox]
\textbf{\textcolor{orbisblue}{\faLightbulb\ Dica:}} A Orbis lê o contexto recente da conversa e adapta o tom. Em grupos agitados responde de forma mais directa; em grupos calmos é mais elaborada.
\end{tcolorbox}

\subsection{Conversa Natural}

Qualquer membro pode conversar livremente com a Orbis:

\begin{itemize}[leftmargin=1.5em, itemsep=2pt]
  \item Fazer perguntas gerais ou sobre o grupo
  \item Pedir opiniões ou sugestões
  \item Pedir para mencionar membros específicos ou todos
\end{itemize}

\textbf{Exemplos:}
\begin{itemize}[leftmargin=1.5em, itemsep=2pt]
  \item \textit{``Orbis, qual é o tema da reunião de hoje?''}
  \item \textit{``@Orbis chama o João para ver isto''}
  \item \textit{``Orbis chama toda a gente''}
\end{itemize}

\subsection{Comandos por Linguagem Natural (Admins)}

Os administradores podem dar ordens à Orbis em linguagem natural, sem usar comandos com \texttt{/}:

\begin{center}
\begin{tabular}{L{6.5cm} L{6.5cm}}
\toprule
\rowcolor{tabhead}\textcolor{white}{\textbf{O que dizes}} & \textcolor{white}{\textbf{O que a Orbis faz}} \\
\midrule
\textit{``Orbis expulsa o Pedro''} & Executa \cmd{/kick @Pedro} \\
\rowcolor{orbisLight}\textit{``Orbis bane o Rui por ameaças''} & Executa \cmd{/ban @Rui} \\
\textit{``Orbis muta o Carlos por 30 minutos''} & Executa \cmd{/mute @Carlos 30m} \\
\rowcolor{orbisLight}\textit{``Orbis avisa o João por spam''} & Executa \cmd{/warn @João} \\
\textit{``Orbis tranca o grupo''} & Executa \cmd{/trancar} \\
\rowcolor{orbisLight}\textit{``Orbis abre o grupo''} & Executa \cmd{/abrir} \\
\textit{``Orbis promove o Miguel''} & Executa \cmd{/promover @Miguel} \\
\rowcolor{orbisLight}\textit{``Orbis anuncia que a reunião é às 18h''} & Envia anúncio formatado \\
\bottomrule
\end{tabular}
\end{center}

\begin{tcolorbox}[warnbox]
\textbf{\textcolor{orbisYellow}{\faExclamationTriangle\ Nota:}} A Orbis só executa acções por linguagem natural se tiver \textbf{alta confiança} na interpretação. Em caso de ambiguidade, não age e pede confirmação.
\end{tcolorbox}

\subsection{Moderação Autónoma}

A Orbis monitoriza todas as mensagens e age automaticamente em caso de violações, \textbf{sem necessidade de comandos}:

\begin{center}
\begin{tabular}{L{3.5cm} C{2.5cm} L{6.5cm}}
\toprule
\rowcolor{tabhead}\textcolor{white}{\textbf{Tipo de Violação}} & \textcolor{white}{\textbf{Severidade}} & \textcolor{white}{\textbf{Acção}} \\
\midrule
Spam, flood, links & Baixa & Apaga mensagem + aviso \\
\rowcolor{orbisLight}Insultos, desrespeito & Média & Aviso ou mute curto \\
Insultos graves, assédio & Alta & Mute ou expulsão \\
\rowcolor{orbisLight}Ameaças, discurso de ódio & Crítica & Expulsão ou ban imediato \\
\bottomrule
\end{tabular}
\end{center}

\textbf{Escalada progressiva por reincidência:}
\[
\text{Aviso} \;\longrightarrow\; \text{Mute} \;\longrightarrow\; \text{Expulsão} \;\longrightarrow\; \text{Ban}
\]

\begin{tcolorbox}[notebox]
\textbf{Administradores nunca são moderados automaticamente.}
\end{tcolorbox}

\subsection{Modo Silêncio}

Quando o modo silêncio está activo, a IA \textbf{não responde a membros} fora do horário configurado. Os administradores continuam a ter acesso em qualquer hora.

\textbf{Configurar:} \cmd{/config silencio} \param{<hora\_início> <hora\_fim>}\\
\textbf{Exemplo:} \cmd{/config silencio 22 6} — silêncio das 22h às 6h


% ═════════════════════════════════════════════════════════════════════════════
\section{Funcionalidades Extra}
\label{sec:extra}
% ═════════════════════════════════════════════════════════════════════════════

\subsection{Sistema de Reputação e Níveis}

O sistema de reputação recompensa membros activos e penaliza infractores.

\textbf{Como funciona:}
\begin{itemize}[leftmargin=1.5em, itemsep=2pt]
  \item Cada mensagem enviada ganha \textbf{+2 pontos}
  \item Cada infracção detectada perde \textbf{-20 pontos}
  \item Ao subir de nível, a Orbis anuncia no grupo
\end{itemize}

\begin{center}
\begin{tabular}{C{1.5cm} L{3cm} C{3cm}}
\toprule
\rowcolor{tabhead}\textcolor{white}{\textbf{Nível}} & \textcolor{white}{\textbf{Nome}} & \textcolor{white}{\textbf{Pontos necessários}} \\
\midrule
1 & Novato & 0 \\
\rowcolor{orbisLight}2 & Membro & 100 \\
3 & Activo & 300 \\
\rowcolor{orbisLight}4 & Veterano & 700 \\
5 & Elite & 1500 \\
\bottomrule
\end{tabular}
\end{center}

\textbf{Activar:} \cmd{/config reputacao on}\\
\textbf{Ver reputação:} \cmd{/reputacao} \param{[@membro]}\\
\textbf{Ver ranking:} \cmd{/top}

\subsection{Resumo Diário Automático}

A Orbis pode enviar automaticamente um resumo da conversa do dia, gerado por IA.

\begin{tcolorbox}[cmdbox, title={\faCalendarCheck\ /resumodiario}]
Activa ou desactiva o resumo diário automático, com hora configurável.\\
\textbf{Uso:} \cmd{/resumodiario} \param{on|off [hora]}\\
\textbf{Exemplo:} \cmd{/resumodiario on 22} — envia resumo todos os dias às 22h
\end{tcolorbox}

\begin{tcolorbox}[cmdbox, title={\faPaperPlane\ /resumoagora}]
Envia o resumo do dia imediatamente, sem esperar pelo horário agendado.\\
\textbf{Uso:} \cmd{/resumoagora}
\end{tcolorbox}

\subsection{Agendamentos Recorrentes}

Permite programar mensagens automáticas que se repetem diariamente ou apenas em dias úteis.

\begin{center}
\begin{tabular}{L{3cm} L{10cm}}
\toprule
\rowcolor{tabhead}\textcolor{white}{\textbf{Tipo}} & \textcolor{white}{\textbf{Descrição}} \\
\midrule
\param{DAILY} & Todos os dias à hora definida \\
\rowcolor{orbisLight}\param{MON-FRI} & Apenas segunda a sexta à hora definida \\
\bottomrule
\end{tabular}
\end{center}

\textbf{Exemplos:}
\begin{itemize}[leftmargin=1.5em, itemsep=2pt]
  \item \cmd{/agendarcorrente DAILY 09:00 Bom dia a todos! ☀️}
  \item \cmd{/agendarcorrente MON-FRI 17:30 Fim do expediente. Bom descanso!}
\end{itemize}

\subsection{Whitelist de Links}

Quando o anti-link está activo, podes permitir links de domínios específicos (ex: YouTube, Google Drive).

\begin{itemize}[leftmargin=1.5em, itemsep=2pt]
  \item \textbf{Adicionar:} \cmd{/config whitelist add youtube.com}
  \item \textbf{Remover:} \cmd{/config whitelist remove youtube.com}
  \item \textbf{Limpar tudo:} \cmd{/config whitelist clear}
\end{itemize}

% ═════════════════════════════════════════════════════════════════════════════
\section{Auto-Moderação — Referência Rápida}
\label{sec:automod}
% ═════════════════════════════════════════════════════════════════════════════

A tabela seguinte resume o comportamento automático do Orbis sem intervenção de admins:

\begin{center}
\begin{longtable}{L{3.5cm} L{3.5cm} L{6cm}}
\toprule
\rowcolor{tabhead}\textcolor{white}{\textbf{Detecção}} & \textcolor{white}{\textbf{Condição}} & \textcolor{white}{\textbf{Acção}} \\
\midrule
\endhead
Link proibido & Anti-link activo & Apaga + aviso \\
\rowcolor{orbisLight}Spam de texto & Anti-spam activo & Apaga + aviso \\
Flood (muitas msgs) & Anti-flood activo & Apaga + aviso \\
\rowcolor{orbisLight}Avisos acumulados & Limite atingido & Expulsão automática \\
Insultos leves & IA activa & Aviso \\
\rowcolor{orbisLight}Insultos graves & IA activa & Mute ou expulsão \\
Ameaças / ódio & IA activa & Ban imediato \\
\rowcolor{orbisLight}Reincidência (2+ strikes) & IA activa & Escalada da punição \\
\bottomrule
\end{longtable}
\end{center}

\begin{tcolorbox}[dangerbox]
\textbf{\textcolor{orbisRed}{\faBan\ Acções irreversíveis:}} \cmd{/ban} e a moderação autónoma por conteúdo crítico resultam em ban permanente. Para reverter, usa \cmd{/adicionar} com o número do membro.
\end{tcolorbox}

% ═════════════════════════════════════════════════════════════════════════════
\section{Referência Rápida de Todos os Comandos}
\label{sec:referencia}
% ═════════════════════════════════════════════════════════════════════════════

\subsection{Comandos para Todos os Membros}

\begin{center}
\begin{tabular}{L{5cm} L{8cm}}
\toprule
\rowcolor{tabhead}\textcolor{white}{\textbf{Comando}} & \textcolor{white}{\textbf{Descrição}} \\
\midrule
\cmd{/menu} & Lista de todos os comandos \\
\rowcolor{orbisLight}\cmd{/info} & Informações do grupo \\
\cmd{/admins} & Lista de administradores \\
\rowcolor{orbisLight}\cmd{/regras} & Regras do grupo \\
\cmd{/stats} & Estatísticas do grupo \\
\rowcolor{orbisLight}\cmd{/sondagem pergunta | op1 | op2} & Criar enquete \\
\cmd{/denunciar} & Denunciar mensagem (responde a ela) \\
\rowcolor{orbisLight}\cmd{/reputacao [@membro]} & Ver reputação \\
\cmd{/top} & Top 10 membros \\
\rowcolor{orbisLight}\cmd{/resumo [N]} & Resumir conversa com IA \\
\cmd{/instrucoes [secção]} & Instruções detalhadas \\
\bottomrule
\end{tabular}
\end{center}

\subsection{Comandos para Administradores}

\begin{center}
\begin{longtable}{L{5.5cm} L{7.5cm}}
\toprule
\rowcolor{tabhead}\textcolor{white}{\textbf{Comando}} & \textcolor{white}{\textbf{Descrição}} \\
\midrule
\endhead
\cmd{/warn @membro [motivo]} & Advertir membro \\
\rowcolor{orbisLight}\cmd{/warnings @membro} & Ver avisos \\
\cmd{/clearwarn @membro} & Limpar avisos \\
\rowcolor{orbisLight}\cmd{/historico @membro} & Histórico de infrações \\
\cmd{/kick @membro} & Expulsar \\
\rowcolor{orbisLight}\cmd{/ban @membro [motivo]} & Banir \\
\cmd{/banidos} & Lista de banidos \\
\rowcolor{orbisLight}\cmd{/adicionar <número>} & Readmitir banido \\
\cmd{/mute @membro [tempo]} & Mutar \\
\rowcolor{orbisLight}\cmd{/unmute @membro} & Desmutar \\
\cmd{/promover @membro} & Promover a admin \\
\rowcolor{orbisLight}\cmd{/despromover @membro} & Despromover \\
\cmd{/trancar} & Só admins enviam \\
\rowcolor{orbisLight}\cmd{/abrir} & Todos podem enviar \\
\cmd{/anunciar <texto>} & Anúncio formatado \\
\rowcolor{orbisLight}\cmd{/deletar} & Apagar mensagem citada \\
\cmd{/agendamento <tempo> <msg>} & Agendar mensagem única \\
\rowcolor{orbisLight}\cmd{/agendarcorrente <tipo> <HH:MM> <msg>} & Agendamento recorrente \\
\cmd{/agendamentos} & Ver agendamentos activos \\
\rowcolor{orbisLight}\cmd{/cancelaragendamento <id>} & Cancelar agendamento \\
\cmd{/denuncias} & Ver denúncias \\
\rowcolor{orbisLight}\cmd{/membros} & Lista de membros \\
\cmd{/resumodiario on|off [hora]} & Resumo diário automático \\
\rowcolor{orbisLight}\cmd{/resumoagora} & Enviar resumo agora \\
\cmd{/config [opção] [valor]} & Ver/alterar configurações \\
\rowcolor{orbisLight}\cmd{/antilink on|off} & Anti-link \\
\cmd{/antispam on|off} & Anti-spam \\
\rowcolor{orbisLight}\cmd{/antiflood on|off} & Anti-flood \\
\cmd{/setregras <texto>} & Definir regras \\
\rowcolor{orbisLight}\cmd{/setwelcome <msg>} & Mensagem de boas-vindas \\
\cmd{/setadmingroup <ID>} & Grupo de notificações \\
\bottomrule
\end{longtable}
\end{center}

% ═════════════════════════════════════════════════════════════════════════════
\section{Boas Práticas}
\label{sec:boas-praticas}
% ═════════════════════════════════════════════════════════════════════════════

\begin{enumerate}[leftmargin=1.5em, itemsep=6pt]
  \item \textbf{Configura o grupo de admins} com \cmd{/setadmingroup} para receberes notificações de moderação em tempo real.
  \item \textbf{Define as regras} com \cmd{/setregras} para que os membros saibam o que é permitido.
  \item \textbf{Activa o anti-link} e adiciona à whitelist apenas os domínios que queres permitir.
  \item \textbf{Usa o modo silêncio} para evitar que a IA responda a horas inapropriadas.
  \item \textbf{Activa a reputação} para incentivar participação saudável e identificar membros problemáticos.
  \item \textbf{Usa \cmd{/historico}} antes de banir — verifica o historial do membro para tomar decisões informadas.
  \item \textbf{Não uses \cmd{/ban} levianamente} — é permanente. Prefere \cmd{/warn} ou \cmd{/kick} para infracções menores.
  \item \textbf{O bot precisa de ser administrador} do grupo para executar acções de moderação (kick, ban, mute, trancar).
\end{enumerate}

% ═════════════════════════════════════════════════════════════════════════════
% RODAPÉ FINAL
% ═════════════════════════════════════════════════════════════════════════════
\vfill
\begin{center}
\textcolor{orbisMid}{\rule{12cm}{0.6pt}}\\[0.5cm]
{\small Desenvolvido por \textbf{Orion Technologies} \textcolor{gray}{·} \url{https://oriontechn.netlify.app}}\\[0.2cm]
{\small \textcolor{gray}{Orbis v1.0 — Manual do Utilizador — \today}}
\end{center}

\end{document}
